\section{Quantum Event Logic}
\label{sec:quantum}

In quantum logic the events are the closed subspaces of a Hilbert
space~\cite{birkhoff-vonneumann1936}. We restrict to rays, the
one-dimensional subspaces, of $\mathbb{R}^d$. A \emph{context} is an
orthogonal basis, a maximal set of pairwise orthogonal rays, whose events
can be observed jointly. A \emph{state} assigns to each ray a number in
$\ENNReal$ such that the numbers on every context sum to $1$. For a unit
vector $\psi$, the \emph{Born rule} $x_v = (\psi \cdot v)^2 / (v \cdot
v)$ gives a state, and mixtures of such states, the density matrices
(positive semidefinite operators of trace $1$), are the \emph{quantum
states}. For $d \ge 3$, Gleason's
theorem~\cite{gleason1957} states that every state on all rays of
$\mathbb{R}^d$ is a quantum state.

This section tests the construction of Section~\ref{sec:recipe} on this
event logic. The mathematics is known: the Kochen--Specker
theorem~\cite{kochen-specker1967}, its $18$-ray
instance~\cite{cabello1996}, and the bounds of the graph-theoretic
approach to contextuality (the impossibility of assigning outcomes
independently of the context in which they are measured)~\cite{cabello2014}. Cox's method has been
applied to non-distributive lattices of this kind, giving
non-Kolmogorovian probability measures for quantum
mechanics~\cite{holik2014}. What is added is its
reading in terms of the construction, with the same graph polytopes as
Section~\ref{sec:second-perfect}, and its mechanization.

The presentation of Definition~\ref{def:mix} takes the rays as events
and, as certain states, the $\{0,1\}$-valued states, which give
exactly one ray of every context the value $1$. By the Kochen--Specker
theorem~\cite{kochen-specker1967}, for $d \ge 3$ there are none. We
formalize a finite instance.

\begin{definition}[Cabello's set\leanref{cabello-set}]
\label{def:cabello}
Cabello, Estebaranz, and Garc\'ia-Alcaine~\cite{cabello1996} give $18$
vectors of $\mathbb{R}^4$ with coordinates in $\{0, \pm 1\}$, and $9$ sets
of $4$ of them:
\[
\begin{array}{l}
\{0001, 0010, 1100, 1\bar{1}00\},\ \{0001, 0100, 1010, 10\bar{1}0\},\
\{1\bar{1}1\bar{1}, 1\bar{1}\bar{1}1, 1100, 0011\},\\
\{1\bar{1}1\bar{1}, 1111, 10\bar{1}0, 010\bar{1}\},\
\{0010, 0100, 1001, 100\bar{1}\},\ \{1\bar{1}\bar{1}1, 1111, 100\bar{1}, 01\bar{1}0\},\\
\{11\bar{1}1, 111\bar{1}, 1\bar{1}00, 0011\},\
\{11\bar{1}1, \bar{1}111, 1010, 010\bar{1}\},\
\{111\bar{1}, \bar{1}111, 1001, 01\bar{1}0\},
\end{array}
\]
where $\bar{1}$ stands for $-1$.
\end{definition}

\begin{theorem}[Kochen--Specker for Cabello's set\leanref{basis-orth}\leanref{no-ks}]
\label{thm:ks}
Each of the $9$ sets is an orthogonal basis of $\mathbb{R}^4$, and each of
the $18$ vectors lies in exactly two of them. No function from the $18$
vectors to $\{0,1\}$ takes the value $1$ at exactly one vector of every
basis.
\end{theorem}

\noindent
The first two statements are checked by evaluation. The third is a
parity count. Summing such a function over the $9$ bases gives $9$, and
counts every vector twice, so the sum is even.

\begin{theorem}[no certain states\leanref{no-certain-state}\leanref{born-not-mix}]
\label{thm:quantum}
For the presentation whose events are the $18$ vectors and whose certain
states are the $\{0,1\}$-valued functions summing to $1$ on every basis,
the set of certain states and $\mathrm{Mix}$ of it are empty. For every
$\psi \in \mathbb{R}^4$ with $\psi \cdot \psi = 1$, the Born-rule
assignment sums to $1$ on every basis, so it is a state that is not in
$\mathrm{Mix}$ of the certain states.
\end{theorem}

\noindent
The Born-rule statement is Parseval's identity
$\sum_{v \in B} (\psi \cdot v)^2 / (v \cdot v) = \psi \cdot \psi$ for
each orthogonal basis $B$, checked basis by basis\leanref{parseval}.

\paragraph{Allowing divergence.}
In a probabilistic coherence space an element may have total mass below
$1$, the missing mass being the probability of
divergence~\cite{danos-ehrhard2011}. Allowing the same here, a state sums to
at most $1$ on every context. The certain states are then the $\{0,1\}$-valued functions
with at most one ray of value $1$ in each basis, the indicators of the
stable sets of the graph that joins two rays when they lie in a common
basis. These exist, and the construction is no longer empty.

\begin{theorem}[a stable-set bound\leanref{quantum-partial}]
\label{thm:quantum-partial}
Every mixture of these certain states has total weight at most $4$ on the
$18$ rays, and $4$ is attained. Every Born-rule assignment of a unit
vector has total weight $\tfrac92$. In particular no Born-rule assignment
is a mixture of these certain states.
\end{theorem}

\noindent
Summing over the nine bases counts every ray twice. A certain state has
weight at most $1$ on each basis, so its total weight is at most
$\tfrac92$, hence at most $4$ as it is an integer. A Born-rule assignment
has weight $1$ on each basis. Summed over the bases, every mixture of
certain states has weight at most $8$, so it diverges with probability at
least $\tfrac19$ on some basis, while a quantum state diverges on none.
The bound $4$ is the independence number of the graph, and the quantum
value $\tfrac92$ does not depend on the state. This is the comparison of
Section~\ref{sec:second-perfect} for a different graph, and the gap
between $4$ and $\tfrac92$ is the Kochen--Specker contradiction in
quantitative form~\cite{cabello2008,cabello2014}.

Quantum probability therefore fails criterion (1) of
Section~\ref{sec:intro}: no state is $\{0,1\}$-valued, so the
$\{0,1\}$-valued states cannot reconstruct the event logic.
Theorem~\ref{thm:quantum} proves this for Cabello's set. For all rays of
$\mathbb{R}^d$ with $d \ge 3$ it follows from the Kochen--Specker theorem, which is not
formalized. With normalized certain states, the construction of
Section~\ref{sec:recipe} produces nothing at this event logic: with no
certain states, every affine law is
valid on the certain states, including $1 \le 0$, and
Theorem~\ref{thm:mix-eq-laws} returns the empty set. The mixtures of
certain states are the noncontextual hidden-variable models of the event
logic, and Theorem~\ref{thm:ks} is the statement that this event logic
has none~\cite{kochen-specker1967}. For lattice valuations,
Theorem~\ref{thm:m3-no-sharp} shows the corresponding absence on the
non-distributive lattice $M_3$, which has no sharp valuation.

\subsection{Local Certainty}
\label{sec:quantum-local}

The construction of Section~\ref{sec:recipe} mixes states that are
certain on every context at once, and the Kochen--Specker theorem rules
these out. A weaker requirement asks for certainty only within each
context.

\begin{definition}[local and global calculi\leanref{local-calc}]
\label{def:local}
Let $E$ be a finite set of outcomes and $\mathcal{C}$ a family of subsets
of $E$, the contexts. The \emph{local calculus} is the set of
$x : E \to \ENNReal$ with $\sum_{e \in C} x_e = 1$ for every context $C$,
and the \emph{global calculus} is $\mathrm{Mix}$ of its $\{0,1\}$-valued
elements. With divergence allowed, the sums are at most $1$.
\end{definition}

\noindent
An element of the local calculus restricts on every context to an element
of that context's classical calculus. Global states are local, and with a
single context the two calculi agree and are classical probability on the
outcomes\leanref{local-calc}, as criterion (2) of Section~\ref{sec:intro}
requires. In the sheaf-theoretic approach to
contextuality~\cite{abramsky-brandenburger2011}, the local calculus
corresponds to the compatible families of local distributions
(distributions on the contexts that agree on their overlaps) and the
global calculus to the global sections (distributions on assignments of
outcomes to all events), the noncontextual hidden-variable models. For
hypergraphs of contexts (families of contexts on a common set of
outcomes), Ac\'in, Fritz, Leverrier, and Sainz
study the same two sets, the general probabilistic models and the
classical models~\cite{acin2015}. Contexts of this kind are
the tests (sets of outcomes of one experiment) of Foulis and
Randall~\cite{foulis-randall1972}.

\begin{theorem}[local certainty in linear logic\leanref{loc-linear}]
\label{thm:loc-linear}
For a formula $A$ of Definition~\ref{def:mall}, take as contexts the
supports of the $\{0,1\}$-valued elements of $P(A^{\perp})$, with
divergence allowed. The calculus of $A$ is contained in $P(A)$, which is
contained in the local calculus, and $P(A)$ equals the local calculus iff
the calculus of $A^{\perp}$ equals $P(A^{\perp})$. At $T$ the first
inclusion is strict and the second is an equality. At $(T \otimes
T)^{\perp}$ the second inclusion is strict.
\end{theorem}

\noindent
The local calculus is the orthogonal of the calculus of $A^{\perp}$, and
$P(A)$ is the orthogonal of $P(A^{\perp})$. Both the calculus and $P$ are
biorthogonally closed (Section~\ref{sec:second-perfect}), which gives the
equivalence. At $T$ the dual is the core, which has no gap
(Theorem~\ref{thm:second-exact}), and the first inclusion is strict by
Theorem~\ref{thm:pentagon}. At $(T \otimes T)^{\perp}$ the indicator of
Shannon's set (Theorem~\ref{thm:shannon}) is a $\{0,1\}$-valued element
of the local calculus outside $P$. So the probabilistic coherence space
is in general neither the global nor the local calculus of its formula.
It is determined by the duality between a formula and its negation. We
did not find this placement in the literature.

\begin{theorem}[local certainty in quantum logic\leanref{loc-quantum}]
\label{thm:loc-quantum}
On Cabello's set with the nine bases of Definition~\ref{def:cabello} as
contexts, the local calculus is the set of states and the global calculus
is empty. Born-rule assignments are local. The function $h$ with value
$\tfrac12$ on the nine rays $\bar{1}111$, $0001$, $0010$, $0100$,
$1\bar{1}\bar{1}1$, $1\bar{1}1\bar{1}$, $11\bar{1}1$, $111\bar{1}$,
$1111$ and $0$ elsewhere is local and is not a mixture of Born-rule
assignments.
\end{theorem}

\noindent
Every basis contains exactly two of the nine rays, so $h$ is local. The
ray $0001$ lies in the span of the orthogonal rays $1001$ and
$100\bar{1}$, and every Born-rule assignment satisfies
$x_{0001} \le x_{1001} + x_{100\bar{1}}$, since the difference of the two
sides is $\psi_1^2$. The function $h$ gives these rays the values
$\tfrac12$, $0$, and $0$. The local calculus sees orthogonality within
contexts but not the order of the subspaces across them, the
monotonicity law of Definition~\ref{def:valuation}. On a finite set of
rays local certainty therefore admits states that are not mixtures of
Born-rule assignments of those rays. This does not decide whether $h$ is
quantum for some other assignment of projections to the rays, in a Hilbert
space of any dimension, with the same contexts. Ac\'in, Fritz, Leverrier,
and Sainz state that they do not know whether every local state on
Cabello's set is quantum in that sense~\cite{acin2015}. On
all rays of $\mathbb{R}^d$ with $d \ge 3$ it does not: there every local
state is a quantum state by Gleason's theorem~\cite{gleason1957}, which
is not formalized. Holik, Plastino, and Sáenz reach the quantum states
from Cox's approach on non-distributive lattices~\cite{holik2014}.

Criterion (1) must be weakened as well, since no state is
$\{0,1\}$-valued. A weaker form asks that the events each state is
certain of, together with the events of probability $0$, determine the
event logic, a condition in the spirit of Mackey's order-determining sets
of states (sets of states whose values determine the order of the event
lattice)~\cite{mackey1963}.

\begin{proposition}[Born states determine the event
logic\leanref{born-ray}]
\label{prop:born-ray}
For rays $u, v$ of Cabello's set, let $x^u$ be the Born-rule assignment
of the unit vector of $u$. Then $x^u_v = 0$ iff $u$ and $v$ are
orthogonal, and $x^u_v = 1$ iff $u = v$.
\end{proposition}

\noindent
So the extension of quantum event logic to a probability calculus needs
two changes to the construction. Certainty is required within each
context and not on all of them at once, and criterion (1) is read through
the certainty and impossibility of events instead of through
$\{0,1\}$-valued states. On all rays of $\mathbb{R}^d$, $d \ge 3$, the
resulting calculus is the set of quantum states. On a finite fragment it
is larger, and excluding the extra states needs the order of the
subspace lattice across contexts.

\section{Four Categorical Readings}
\label{sec:categorical}

The construction $S \mapsto \mathrm{Mix}(S)$ and criteria (1)--(3) can be
read in four categorical settings: free algebras of a probability monad,
orthogonality, effect algebras, and limits. The mathematics of each
reading is standard, and we credit it below. This section records what
each reading says about the calculi of this paper, with the new
statements mechanized.

\subsection{Free Convex Algebras and Simplices}
\label{sec:cat-simplex}

Convex sets are the algebras of the finite distribution monad
$D$~\cite{swirszcz1974,fritz2009,jacobs2011}. The monad sends a set to
the finitely supported probability distributions on it, and an algebra
of it is a set with an operation that evaluates such a distribution as a
mixture. For a presentation $(E, S)$, the inclusion
$S \subseteq \ENNReal^{E}$ extends to a unique affine map (a map
preserving mixtures) $D(S) \to \ENNReal^{E}$, and its image is
$\mathrm{Mix}(S)$. The certain states are the image of the unit of $D$
(the map sending an element to the point mass at it). Abramsky and Brandenburger
describe the noncontextual models of a measurement scenario in the same
way, as the image of the distributions on global
assignments~\cite{abramsky-brandenburger2011}, and Pitowsky's
correlation polytopes are the same sets~\cite{pitowsky1991}. In these
terms, criterion (1) with closure under mixtures says that the calculus
contains the image of the free algebra on the certain states, and
criterion (3) says that it lies in the solution set of their laws.
Theorem~\ref{thm:unique} says the two coincide.

The calculus is a \emph{simplex} when the map $D(S) \to \mathrm{Mix}(S)$
is injective, that is, when every state is a mixture of certain states in
only one way. For finite $S$ this holds iff $S$ is affinely independent.
In the theory of generalized probabilistic theories, which describes a
system by a convex set of states and a set of effects (the affine maps
from states to $[0,1]$ that give the probabilities of outcomes), a system
is called classical when its state space is a simplex~\cite{barrett2007,plavala2023}.

\begin{theorem}[simplices\leanref{mix-deltapoint-inj}\leanref{lin-not-simplex}\leanref{rep-unique}]
\label{thm:cat-simplex}
(i) For a finite poset $P$, two mixtures of point valuations on
$\mathrm{Low}(P)$ that agree on every lower set have the same weights.
(ii) At $\mathbf{2} \multimap \mathbf{2}$ the even mixture of the
identity and negation equals the even mixture of the two constant
programs.
(iii) Let $(E, S)$ be a presentation with an event certain in every
$s \in S$, and for all events $a, b$ an event certain in $s \in S$ iff
$a$ and $b$ are. Two probability measures on $\{0,1\}^{E}$ that vanish
on $\mathrm{bad}_F$ for every window $F$ and give each event the same
mass agree on the $\sigma$-algebra generated by the events. The sharp
valuations of every frame satisfy the hypothesis, with $\top$ and
$a \sqcap b$.
\end{theorem}

\noindent
Part (i) evaluates the mixtures at $\downarrow\! p$ and at
$\downarrow\! p \setminus \{p\}$, whose difference is the weight of $p$.
It is the Möbius inversion (the inversion of summation over a partial
order) of a valuation on a finite distributive lattice~\cite{rota1964,geissinger1973}. Part (ii) is the centre of the
square of $2 \times 2$ stochastic matrices, Barrett's gbit~\cite{barrett2007}.
Part (iii) applies the uniqueness of measures on a $\pi$-system (a
family of sets closed under finite intersections): the
finite intersections of events form one, and off the bad set of a window
an intersection of events is the event of their meet.

So the valuation calculus of a finite frame is a simplex. The calculus
of linear logic is not one already at first order. Part (iii) with
Theorem~\ref{thm:barycenter} and Corollary~\ref{cor:frame-unique} gives,
in a classical metatheory, that every valuation on a frame $\Omega$ is the
barycenter of a probability measure on $\{0,1\}^{\Omega}$ that vanishes
on every bad set, unique on the $\sigma$-algebra of the events. The certain states here are all sharp
valuations, the prime filters, and not only the points: on the powerset
of $\mathbb{N}$ the non-principal ultrafilters can carry mass. With the
extension theorem of Smiley, Horn and Tarski, and
Pettis~\cite{smiley1944,horn-tarski1948,pettis1951} and the
correspondence between additive functionals on a Boolean algebra and
Radon (inner regular) measures on its Stone space~\cite{fremlin2003},
this makes the valuations a Bauer simplex (a compact convex set whose
extreme points form a closed set and in which every point is the
barycenter of a unique probability measure on them). For Boolean algebras, and more generally for
MV-algebras, this is Kroupa's theorem~\cite{kroupa2006}, and for
continuous valuations on stably compact spaces (the spaces of domain
theory whose patch topology is compact Hausdorff) it is the theorem of
Lawson and of Alvarez-Manilla, Jung, and Keimel~\cite{ajk2004}. For
valuations on frames without continuity we found it only as a
combination of these parts. It is not formalized.

This gives a reading of intuitionistic probability as a generalized
probabilistic theory, which we did not find in the literature. On a
finite frame its state space is a simplex, so it is classical in that
sense. What is not classical is the set of effects, the convex hull of
the indicators of opens, which is closed under pointwise $\wedge$ and
$\vee$ and not under complement. Restricted theories in the sense of
Janotta and Lal assume closure under complement~\cite{janotta-lal2013}.
Filippov, Gudder, Heinosaari, and Leppäjärvi derive it from the
requirement that every effect belong to a measurement, a family of
effects summing to the unit effect~\cite{filippov2020}. Here the
indicator of an open $U$ belongs to a measurement iff $U$ is clopen,
since the other effects of the measurement sum to the indicator of the
complement of $U$, and a $\{0,1\}$-valued nonnegative combination of
indicators of opens is the indicator of their union. The
effects do satisfy the conditions that Plávala lists: they are convex,
contain $0$ and the unit, and separate the states~\cite{plavala2023}.
With the states taken as the measures on certain states, part (iii) says
that the opens separate them. The states cannot be recovered from the
effects by the duality of generalized probabilistic theories. On the
Sierpiński space, with points $\bot \le \top$ and the open $\{\top\}$,
the functional $(1 + t)\,\delta_{\top} - t\,\delta_{\bot}$ is nonnegative on
every effect and $1$ at the unit for every $t \ge 0$, and gives
$\{\top\}$ the value $1 + t$. The calculus fixes its states by mixing
certain states instead.
The simplex reading depends on enough certain
states. Heunen, Landsman, and Spitters represent the quasi-states
(functionals linear on each commuting family of observables) of a
quantum system, which are its states when Gleason's theorem applies, as
valuations on a locale internal to a topos (a category that behaves as a
universe of sets with an intuitionistic internal logic), and this locale has no
points when the Hilbert space has dimension greater than
$2$~\cite{heunen-landsman-spitters2009}. So quantum states are
intuitionistic valuations without certain states, and the separation of
intuitionistic from quantum probability above holds in a classical
metatheory, where the prime ideal theorem supplies the certain states.

\subsection{Orthogonality}
\label{sec:cat-orth}

States and affine laws form a Galois connection (a pair of
order-reversing maps, sending a set of states to the laws they all
satisfy and a set of laws to the states satisfying all of them, whose
composite is a closure operator), and Theorem~\ref{thm:mix-eq-laws} says
that $\mathrm{Mix}(S)$ is the closure of $S$. A probabilistic coherence space is closed for a smaller class of
tests, the nonnegative $y$ with $\sum_w x_w y_w \le 1$, by Girard's
bipolar theorem~\cite{girard2004}, and the category of probabilistic
coherence spaces is an orthogonality category in the sense of Hyland and
Schalk~\cite{hyland-schalk2003,fiore-galal-jafarrahmani2023,tsukada-asada2022}.
An object of such a category is a tight pair, a set of states and a set
of tests each of which is the orthogonal of the other. We ask when the
calculi of a formula and of its negation form a tight pair.

\begin{theorem}[duality\leanref{orth-recipe}\leanref{tight}]
\label{thm:cat-duality}
For every formula $A$ of Definition~\ref{def:mall}, the orthogonal of
the calculus of $A$ is the local calculus of $A^{\perp}$, and the
orthogonal of the local calculus of $A$ is the calculus of $A^{\perp}$.
The calculi of $A$ and of $A^{\perp}$ form a tight pair iff the calculus
of $A$ equals its local calculus, iff neither $A$ nor $A^{\perp}$ has a
gap. At the core $T^{\perp}$ they do not form a tight pair.
\end{theorem}

\noindent
The first two statements combine the biorthogonal closure of the
calculus with the fact that orthogonals do not see mixing. The third
follows with Theorem~\ref{thm:loc-linear}, and the last with
Theorem~\ref{thm:second-exact} and the gap at $T$
(Theorem~\ref{thm:pentagon}). Linear negation therefore exchanges the
global and the local calculus, and the probabilistic coherence space
$P(A)$, with $P(A^{\perp}) = P(A)^{\perp}$, is the self-dual set between
them. The calculus of certain states is not a model of linear negation,
since it is not tight at $T^{\perp}$. The failure arises at a tensor of
tight types.

\begin{corollary}[tightness and the tensor\leanref{tight-tensor}]
\label{cor:cat-tensor}
For all $n$ and $m$, the local calculus of $\mathbf{n} \multimap
\mathbf{m}$ is $P(\mathbf{n} \multimap \mathbf{m})$, so
$\mathbf{n} \multimap \mathbf{m}$ is tight. For
$n, m, n', m' \ge 2$, the tensor $(\mathbf{n} \multimap \mathbf{m})
\otimes (\mathbf{n'} \multimap \mathbf{m'})$ of two tight types is not
tight.
\end{corollary}

\noindent
The indicator of a row $\{a\} \times \{1, \dots, m\}$ is a certain state
of the dual of $\mathbf{n} \multimap \mathbf{m}$, and these indicators
cut out the row-substochastic matrices, which form
$P(\mathbf{n} \multimap \mathbf{m})$. With
Theorem~\ref{thm:pcs-recipe} the calculus is the local calculus. The
second statement is Theorem~\ref{thm:cat-duality} with the gap of
Theorem~\ref{thm:pentagon}. So the assignment of the certain-state
calculus to formulas is not compositional: tightness at $A$ and at $B$
does not imply tightness at $A \otimes B$. The probabilistic coherence space is compositional by
construction.

The convex geometry is the antiblocking duality of Fulkerson and
Lovász~\cite{fulkerson1971,lovasz1972}. When the certain states of $A$
and of $A^{\perp}$ are the cliques of their coherence graphs, the calculus and the local calculus are the
stable set polytope and the clique-constrained polytope of the
complementary graph, the duality is Knuth's
$\mathrm{STAB}(\overline{G})^{\perp} = \mathrm{QSTAB}(G)$, and
tightness is $\mathrm{STAB} = \mathrm{QSTAB}$, which holds iff the graph is
perfect~\cite{knuth1994,chvatal1975}. Lovász's theta body lies between
the two polytopes and is also self-dual~\cite{knuth1994}, and the chain
of the three is Cabello, Severini, and Winter's classical, quantum, and
exclusivity bounds~\cite{cabello2014}. The set $P(A)$ is not the theta
body: at $T$ it is the local calculus (Theorem~\ref{thm:loc-linear}).
Two sources state parts of this. In an exercise sheet, Ehrhard
assigns to a coherence space $E$ the probabilistic coherence space
$p(E)$ of the $x \ge 0$ with total weight at most $1$ on every clique of
$E^{\perp}$, and asks for proofs that $p(E^{\perp}) = p(E)^{\perp}$ for the
coherence spaces generated from $\mathbf{1}$ by $\perp$ and $\mathbin{\&}$,
and that this fails whenever the five-cycle embeds in
$E$~\cite{ehrhard-td2}. When the certain states of the dual are all its
cliques, $p(E)$ is the local calculus. Borodulin-Nadzieja, Farkas, and
Pelczar-Barwacz show that a family of finite sets and its orthogonal
family are geometrically dual (each family's norm ball is the dual of
the other's) iff they are the cliques and the anticliques (stable sets)
of a perfect graph~\cite{bfp2026}. We did not find the
tight-pair formulation for the calculus of the certain states of a
formula, whose certain states can be fewer than the cliques
(Theorem~\ref{thm:shannon}), its equivalence with the absence of a gap
at the formula and at its negation, or the failure at the core.

At the core the failure of tightness is Bell's theorem. A point
$((x, a), (y, b))$ of the web of $T^{\perp} = (\mathbf{2} \multimap
\mathbf{2}) \otimes (\mathbf{2} \multimap \mathbf{2})$ is the event that
two parties with inputs $x$ and $y$ give outputs $a$ and $b$. Two events
are \emph{locally orthogonal} if some party has equal inputs and
different outputs. This is the opposite of the exclusivity of points of
$T$ in Section~\ref{sec:second-seq}, since the web of $T$ is the web of
$T^{\perp}$ and its points are read as tests instead of outcomes.

\begin{theorem}[Bell at the core\leanref{bell}]
\label{thm:cat-bell}
The $\{0,1\}$-valued elements of $P(T)$ are indicators of sets of
pairwise locally orthogonal events. The Popescu--Rohrlich box, with value
$\tfrac12$ on the events with $a + b \equiv xy \pmod 2$ and $0$
elsewhere, lies in the local calculus of $T^{\perp}$ and not in
$P(T^{\perp})$.
\end{theorem}

\noindent
Two events that are not locally orthogonal are both hit by one
deterministic product strategy (a pair of deterministic programs, one
for each party), which lies in $P(T^{\perp})$ and pairs to $2$ with
their indicator. So the local calculus of $T^{\perp}$ is the set of
functions with total weight at most $1$ on every set of pairwise
locally orthogonal events, the principle of local orthogonality, and no
three events in the support of the box are pairwise locally orthogonal. The box is
outside $P(T^{\perp})$ because every product of two substochastic
matrices satisfies the CHSH inequality, total weight at most $3$ on the
eight winning events, which the box violates with $4$. By
Theorem~\ref{thm:second-exact}, $P(T^{\perp})$ is the calculus, the
mixtures of the cliques of the core, which are the subsets of the
supports of deterministic product strategies. Its normalized elements,
with total weight $1$ for each pair of inputs, are therefore the
Bell-local behaviours. For two
parties local orthogonality is equivalent to no-signalling on normalized
behaviours~\cite{fritz-lo2013}, so on normalized behaviours the local
calculus of $T^{\perp}$ is the no-signalling polytope (the behaviours
in which each party's output distribution does not depend on the other
party's input), which is also Kissinger and Uijlen's tensor of the two
first-order types in their category of causal processes (processes
compatible with a causal order)~\cite{kissinger-uijlen2019}. The identification
of the normalized faces is not formalized. By
Theorem~\ref{thm:loc-linear}, the gap between $P$ and the local calculus at
$T^{\perp}$ and the gap between the calculus and $P$ at $T$ are equivalent,
so the pentagon point at $T$ and the box at $T^{\perp}$ witness one
phenomenon from its two sides.

\subsection{Effect Algebras and Extension of Scalars}
\label{sec:cat-effect}

An effect algebra is a set with a partial commutative associative
addition $\oplus$, a unit $1$, and for each $a$ a unique orthosupplement
$a^{\perp}$ with $a \oplus a^{\perp} = 1$, in which $a \oplus 1$ is
defined only for $a = 0$~\cite{foulis-bennett1994}. Its
states are the maps to $[0,1]$ that preserve $\oplus$ and $1$, so they
satisfy the complement rule. Jacobs and Mandemaker extend a logic to
probability by tensoring its effect algebra with the unit interval, and
they state Gleason's theorem as the equation of the effects on a Hilbert
space with the free effect module (the effect algebra with scalars from
$[0,1]$ freely generated by a given one) on its
projections~\cite{jacobs-mandemaker2012,jacobs2015}. This reading does not reach intuitionistic logic,
for three reasons.

\begin{theorem}[effect-algebra states on frames\leanref{em-of-complemented}\leanref{partial-sum}]
\label{thm:cat-effect}
(i) If every element of a frame has a complement, the frame is Boolean.
(ii) Call $q : \Omega \to \ENNReal$ a partial-sum state if it is
normalized, monotone, and additive on joins of disjoint elements. Every
valuation is a partial-sum state. On the five-element frame of
Example~\ref{thm:no-functional}, some partial-sum state violates
modularity, an affine law valid on every valuation.
\end{theorem}

\noindent
Part (i) is known in substance: in a Heyting effect algebra (an effect
algebra that is also a Heyting algebra in a compatible way) the Heyting
negation lies below the orthosupplement, with equality exactly at the
central elements (those that split the algebra as a
product)~\cite{foulis2000}. Part (ii) uses Darst's
example~\cite{darst1970}. Reading a frame as an effect algebra, with
disjoint join as partial sum and with an orthosupplement, forces it to be
Boolean by (i). Keeping the partial sum and dropping the orthosupplement
gives the partial-sum states, which can violate criterion (3), by (ii)
and Theorem~\ref{thm:unique}. The extension of scalars that does satisfy the
criteria on frames replaces partial addition by modularity. It is the
valuation monoid of Definition~\ref{def:vmonoid}, the monoid of Tarski
and of Horn and Tarski~\cite{tarski1938,horn-tarski1948}, whose
normalized monotone homomorphisms to $\ENNReal$ are the valuations
(Theorem~\ref{thm:vmonoid}). State notions on residuated lattices
(lattices with a monoid operation and implications adjoint to it) that
keep the complement rule fail criterion (1) instead. Bosbach and Riečan
states (two notions of state for residuated lattices, defined by
additivity conditions on their operations) on a Heyting algebra satisfy $s(\neg x) = 1 - s(x)$ and factor
through its regular elements~\cite{dvurecenskij-rachunek2006}. By the
second item of the paragraph on conflation in
Section~\ref{sec:recipe-unique}, on a frame that is not Boolean some sharp
valuation violates this equation, so these states do not contain the
certain states. For quantum logic, the local calculus of
Section~\ref{sec:quantum-local} is the set of states of a test
space~\cite{foulis-randall1972,acin2015}, and on all rays of
$\mathbb{R}^d$ with $d \ge 3$ Gleason's theorem identifies it with the
quantum states, the case that Jacobs and Mandemaker treat.

\subsection{Limits of Finite Calculi}
\label{sec:cat-limit}

For arbitrarily many events, the calculi on finite windows form a
functor (an assignment of sets to windows and of maps to inclusions that
preserves identities and composition) from the opposite of the poset of
finite windows to sets, sending a window $F$ to $\mathrm{Mix}$ of the
restrictions of $S$ to $F$ and an inclusion of windows to restriction.
Its limit is the set of families of window states that agree under
restriction.

\begin{theorem}[limit\leanref{window-limit}]
\label{thm:cat-limit}
For every presentation $(E, S)$, the set of functions satisfying every
finite law valid on $S$, which is the calculus of
Theorem~\ref{thm:unique-closed}, is in bijection with the limit of this
functor, by restriction to windows.
\end{theorem}

\noindent
The proof identifies the limit with the compatible families of window
states and reads a family off its values on singleton windows.
Theorem~\ref{thm:barycenter} gives the rest. So the closure condition
of Theorem~\ref{thm:unique-closed} says that the calculus on all events
is determined by its finite windows. For frames the same holds with
finite sublattices in place of windows: the valuations on a distributive
lattice $L$ are the limit, over its finite sublattices $L_0$, of the
simplices $D(\mathrm{pt}\,L_0)$ of distributions on their prime filters.
Distributive lattices are the filtered colimits (directed unions) of
their finite sublattices, and valuations turn these colimits into
limits, so the valuation functor is the right Kan extension of
$D \circ \mathrm{pt}$ (its extension from finite lattices to all
lattices by taking limits) along the inclusion of finite distributive
lattices~\cite{adamek-rosicky1994}. This is folklore. We did not find it
stated, but Gehrke, Jakl, and Reggio prove the same continuity for
lattice-valued measures on distributive
lattices~\cite{gehrke-jakl-reggio2022}. It is not formalized.

The closest results are codensity theorems. The codensity monad of a
functor $G$ is the right Kan extension of $G$ along itself. It
generalizes the monad $G \circ F$ that $G$ induces when it has a left
adjoint $F$. Leinster shows that the
ultrafilter monad (sending a set to its ultrafilters) is the codensity
monad of the inclusion of finite sets into sets~\cite{leinster2013}. Avery shows that the finitely additive
Giry monad is a codensity monad~\cite{avery2016}, and Van Belle that it
is the codensity monad of the functor sending a finite set to its
simplex of distributions~\cite{vanbelle2022}. The ultrafilters on a set
are the certain states of its powerset frame, and the finitely additive
probabilities are its calculus (Corollary~\ref{cor:frame-unique}). So
Leinster's theorem concerns the certain states of the powerset frame and
Van Belle's its calculus. The valuation functor on distributive lattices
goes from lattices to convex sets and is not an endofunctor (a functor
from a category to itself), which is why
the statement above is a Kan extension and not a codensity monad.
We do not know whether the valuations on frames arise as a codensity
monad in a suitable category of spaces.
